% Created 2026-06-19 Fri 07:54
% Intended LaTeX compiler: pdflatex
\documentclass{article}
\usepackage[utf8]{inputenc}
\usepackage[T1]{fontenc}
\usepackage{graphicx}
\usepackage{longtable}
\usepackage{wrapfig}
\usepackage{rotating}
\usepackage[normalem]{ulem}
\usepackage{amsmath}
\usepackage{amssymb}
\usepackage{capt-of}
\usepackage{hyperref}

\usepackage{fancyhdr}
\usepackage[top=25mm, left=25mm, right=25mm]{geometry}
\usepackage{longtable}
\fancyhead[R]{}
\setlength\headheight{43.0pt}
\usepackage[T1]{fontenc}
\usepackage[utf8]{inputenc}
\usepackage[spanish, ]{babel}
\usepackage[backend=biber]{biblatex}
\addbibresource{/home/emilioalbanfs/BusesDelSistema_G1-G2/FormatoTareas/bibliography.bib}
\author{G1 [Emilio Alban, Ariel Montufar, Karen Suarez, Samuel Churaco]}
\date{2026-06-19}
\title{Toma de Notas - S8 Buses del Sistema}
\hypersetup{
 pdfauthor={G1 [Emilio Alban, Ariel Montufar, Karen Suarez, Samuel Churaco]},
 pdftitle={Toma de Notas - S8 Buses del Sistema},
 pdfkeywords={},
 pdfsubject={},
 pdfcreator={Emacs 27.1 (Org mode 9.7.5)}, 
 pdflang={Español}}
\begin{document}

\maketitle
\tableofcontents

\fancyhead[C]{\includegraphics[scale=0.05]{./logoEPN.jpg}\\
ESCUELA POLITÉCNICA NACIONAL\\FACULTAD DE INGENIERÍA DE SISTEMAS\\
ARQUITECTURA DE COMPUTADORES}
\thispagestyle{fancy}


\section{Estructuras de Interconexión}
\label{sec:org3f28678}

\subsection{Preguntas y Términos clave (Cue):}
\label{sec:org68b66ec}
\begin{description}
\item[{¿Qué es una estructura de interconexión?}] Es el conjunto de líneas
o caminos que conectan los módulos principales (CPU, Memoria, E/S)
para permitir la comunicación entre ellos.
\item[{¿Cuáles son los tres módulos elementales del sistema?}] El
Procesador (CPU), la Memoria y el módulo de Entrada/Salida (E/S).
\item[{¿Qué señales intercambia típicamente la CPU?}] Lee instrucciones y
datos, escribe datos procesados, envía señales de control y recibe
señales de interrupción.
\item[{¿En qué consiste el Acceso Directo a Memoria (DMA)?}] Es un
intercambio directo de datos entre el módulo de E/S y la memoria sin
que la información pase por el procesador, lo que ahorra tiempo de
ejecución.
\end{description}

\subsection{Notas (Notes):}
\label{sec:org18d603c}

\subsubsection{El computador como una red de módulos}
\label{sec:orgaa6b981}
\begin{itemize}
\item Un computador no es una sola pieza, sino una red de módulos que se
comunican constantemente.
\item La estructura de interconexión define los caminos por donde viaja la
información entre estos componentes.
\end{itemize}

\subsubsection{Los Módulos y sus Señales (Basado en Figura 3.15 de Stallings)}
\label{sec:org928b7e6}
\begin{itemize}
\item CPU (El Cerebro): Actúa como el controlador principal.
\begin{itemize}
\item Salidas: Direcciones y señales de control.
\item Entradas: Datos, instrucciones y señales de interrupción.
\end{itemize}
\item Memoria (El Almacén): Constituida por N palabras de igual longitud,
cada una con una dirección única. 
\begin{itemize}
\item Operaciones principales: Lectura y Escritura.
\end{itemize}
\item Módulo de E/S (La Interfaz): Similar a la memoria pero maneja
dispositivos externos y puertos. 
\begin{itemize}
\item Puede enviar señales de interrupción a la CPU para avisar que
terminó una tarea.
\end{itemize}
\end{itemize}

\subsubsection{Tipos de Transferencias de Datos}
\label{sec:orgce3a666}
La estructura debe soportar 5 intercambios fundamentales:
\begin{enumerate}
\item Memoria al Procesador: La CPU lee una instrucción o dato desde la
memoria.
\item Procesador a Memoria: La CPU escribe una unidad de datos en una
dirección específica.
\item E/S al Procesador: La CPU capta datos de un dispositivo externo a
través del módulo de E/S.
\item Procesador a E/S: El procesador envía datos o comandos hacia un
periférico.
\item Memoria a E/S y viceversa (DMA): Transferencia directa de bloques
de datos entre memoria y periféricos sin intervención de la CPU.
\end{enumerate}

\subsection{Resumen (Summary):}
\label{sec:orgf6a8f91}
El computador funciona como una red de tres módulos (CPU, Memoria y
E/S) unidos por una estructura de interconexión. Esta estructura
permite mover datos e instrucciones mediante cinco tipos de
transferencias clave, destacando el DMA por su capacidad de realizar
intercambios directos entre memoria y periféricos para mejorar la
eficiencia del sistema. 


\section{Interconexión punto a punto}
\label{sec:org465b6e7}
\subsection{Preguntas y Términos clave (Cue):}
\label{sec:orgd5255ac}
\begin{itemize}
\item ¿Por qué los buses compartidos tradicionales están siendo
reemplazados en los sistemas modernos? :: Debido a las restricciones
eléctricas al aumentar la frecuencia y la dificultad de realizar
funciones de sincronización y arbitraje en tiempos cada vez más
cortos .
\item[{QPI (QuickPath Interconnect)}] Es una tecnología de interconexión
punto a punto de alto rendimiento introducida por Intel en 2008 para
superar los cuellos de botella de los buses compartidos.
\item ¿Cuáles son las características principales de las arquitecturas
punto a punto? :: Conexiones directas por parejas, arquitectura de
protocolos en capas y transferencia de datos mediante paquetes con
cabeceras y control de errores.
\item[{¿Para qué sirve el carril de reloj en QPI?}] Sirve para sincronizar
el envío de señales, permitiendo que el receptor identifique
exactamente cuándo se transmite cada bit en los 20 carriles de
datos.
\item[{¿Puede haber pérdida de bits? ¿En qué casos?}] Sí, el ruido
eléctrico puede alterar bits. La capa de enlace detecta esto
mediante el código CRC y solicita la retransmisión del paquete
dañado.
\end{itemize}

\subsection{Notas (Notes):}
\label{sec:org00d3ca1}
\subsubsection{El fin del bus compartido}
\label{sec:org8d22d38}
\begin{itemize}
\item El bus compartido fue el estándar durante décadas, pero está siendo
abandonado debido a restricciones eléctricas en sistemas de alta
velocidad.
\item A medida que se incrementa la frecuencia de los buses síncronos
anchos, resulta casi imposible sincronizar y arbitrar los datos a
tiempo.
\item La llegada de los chips multinúcleo multiplicó estas dificultades,
requiriendo menor latencia y mayor tasa de datos para no frenar a
los procesadores.
\end{itemize}
\subsubsection{Características de la Interconexión Punto a Punto}
\label{sec:org4f61726}
\begin{itemize}
\item Este modelo ofrece menor latencia, mayor tasa de datos y mejor
escalabilidad que los buses tradicionales.
\item Presenta 3 características fundamentales:
\begin{enumerate}
\item \textbf{Múltiples conexiones directas:} Enlaces directos por parejas
entre componentes, eliminando la necesidad de arbitraje por el
uso del medio.
\item \textbf{Arquitectura de protocolos en capas:} Estructura similar a redes
(como TCP/IP), organizada en 4 niveles técnicos.
\item \textbf{Transferencia por paquetes:} Los datos se encapsulan en paquetes
con etiquetas de control y códigos de corrección de errores
(CRC), en lugar de enviarse como un flujo de bits en bruto .
\end{enumerate}
\end{itemize}

\subsubsection{El modelo QPI (QuickPath Interconnect)}
\label{sec:org4cd4940}
\begin{itemize}
\item Arquitectura estructurada en 4 capas: Física (señales eléctricas), Enlace (fiabilidad), Encaminamiento (tablas de ruta) y Protocolo (coherencia de caché).
\item \textbf{Unidades de transferencia:} 
\begin{itemize}
\item Capa Física: La unidad es el \textbf{Phit} (20 bits).
\item Capa de Enlace: La unidad es el \textbf{Flit} (80 bits: 72 de mensaje y 8 de CRC).
\end{itemize}
\item \textbf{Transmisión y Hardware:} Utiliza 20 carriles de datos por dirección y señalización diferencial de bajo voltaje (LVDS) para mitigar el ruido.
\item \textbf{Control de Flujo:} Usa un sistema de créditos; el receptor devuelve créditos al emisor cuando sus buffers tienen espacio libre, evitando la saturación del sistema.
\item \textbf{Cálculos de capacidad de transferencia:} 
A velocidades típicas de 6.4 GT/s (transferencias por segundo), enviando 20 bits por Phit en cada transferencia, la capacidad por dirección es:
\(6.4 \text{ GT/s} \times 20 \text{ bits} = 16 \text{ GB/s}\).
Como los enlaces QPI son pares bidireccionales dedicados, la capacidad total de un enlace es:
\(16 \text{ GB/s} \times 2 = 32 \text{ GB/s}\).
\end{itemize}
\subsection{Resumen (Summary):}
\label{sec:orga1c070b}
La evolución tecnológica abandonó los buses compartidos porque las
altas frecuencias eléctricas hacían casi imposible sincronizar y
arbitrar los datos a tiempo. Al evolucionar a arquitecturas de
interconexión punto a punto de alto rendimiento como QPI, se mejora
drásticamente la escalabilidad y la latencia del sistema. Esto se
logra empleando conexiones directas que eliminan el arbitraje, y
gestionando la información estructurada en paquetes a través de
protocolos de capas (Física, Enlace, Enrutamiento y Protocolo)


\section{{\bfseries\sffamily TODO} Mis Tareas [100\%]}
\label{sec:org24bd83b}
\begin{itemize}
\item[{$\boxtimes$}] Leer este documento
\item[{$\boxtimes$}] Escribir el resumen  de la exposición 1 - G1
\end{itemize}

\printbibliography
\end{document}
